% ============================================================
% Digital Discovery Conclusions — honest benchmark (repositioned 2026-06-23)
% ============================================================

\section{Conclusions}\label{sec:conclusions}

We benchmarked how well cheap structure-based machine learning predicts the
experimental singlet--triplet gap $\Delta E_\text{ST}$ of TADF emitters, with every
number reproducible from the deposited code. Three findings stand out.

\paragraph{Structure alone gives a modest, well-characterised predictor.}
Random forests built from the 2D structure alone---Morgan fingerprints
(MAE = \qty{0.091}{\electronvolt}) and RDKit descriptors
(\qty{0.100}{\electronvolt})---predict experimental $\Delta E_\text{ST}$ as well
as NTO descriptors derived from a semi-empirical excited-state calculation
(\qty{0.096}{\electronvolt}), over \num{231} molecules and \num{212} scaffolds, and
the result holds across scaffold, cluster and random cross-validation. All reach only
$R^2 \approx 0.25$--$0.31$ with confidence intervals that include zero, and weak
ranking ($\rho \approx 0.26$--$0.36$); the signal falls below a label-permutation null
($p = 0.001$) yet stands only modestly above a mean-value baseline. The excited-state
calculation is unnecessary for the gap; SHAP on the NTO variant nonetheless ties the
signal to hole--electron overlap and charge-transfer character, keeping it interpretable.

\paragraph{Data, not model capacity, sets the accuracy ceiling.}
Inter-laboratory disagreement (\qty{0.072}{\electronvolt} RMS, up to
\qty{0.56}{\electronvolt}) and the
functional dependence of the computed reference (mean \qty{0.30}{\electronvolt}, up to
\qty{0.58}{\electronvolt} between B3LYP and CAM-B3LYP) bound achievable accuracy near
\qty{0.1}{\electronvolt}; no higher-capacity learner beats the random forest, so the
limit is the data, not the model. Claims of near-DFT precision against such references
are not physically meaningful.

\paragraph{A cautionary methodological note.}
Secondary to the benchmark, regressing $\Delta E_\text{ST}$ onto a feature set that
contains the $S_1$ and $T_1$ excitation energies reproduces the target by construction
($\Delta E_\text{ST}\equiv E_{S_1}-E_{T_1}$); the apparent sub-\num{0.05}~eV accuracy of
such models reflects arithmetic, not learning. This general leakage failure mode---any
target that is an algebraic combination of features will be ``predicted'' by
arithmetic, invisibly to cross-validation---is a reminder that benchmarks must target
the measured gap with features that exclude the constituent energies.

\paragraph{Practical scope.}
From the 2D structure alone, with no quantum chemistry, the model acts as a coarse
triage filter---enriching low-gap candidates \numrange{1.2}{1.4}$\times$ over random
selection, for which we deposit a conformal-bounded shortlist of the enumerated
library---useful for prioritisation but far from a selector, and not a quantitative
photophysical method. Matching a specific lighting niche (horticultural or
human-centric) needs the emission wavelength and oscillator strength, which this
scalar-gap model does not predict, and is left to downstream work.
We make no device-level performance claims: the spin--orbit-coupling-based
$k_\text{RISC}$ and external-quantum-efficiency estimates explored during this work
were molecule-independent and have been withdrawn.

\paragraph{Outlook.}
Progress beyond the present ceiling requires a curated, single-source,
solvent-controlled experimental benchmark and uncertainty-aware models that report
calibrated confidence with each prediction. All data, trained models, and the
scripts that regenerate every number and figure in this work are openly available
(Zenodo DOI: 10.5281/zenodo.17436069), so that the benchmark can be extended and
the limitations we report can be tested directly.
